\icode{E-MD-SI-L3-261114-001}
“血沉”是指红细胞在一定条件下沉降的速度，在医学中具有重要意义。测量“血沉”可将经过处理后的血液放进血沉管内，由于重力作用，血液中的红细胞将会下沉。设血沉管竖直放置且足够深，红细胞的形状为球体。已知红细胞下落受到血液的粘滞阻力表达式为 \(f = 6\pi \eta r v\)，其中 \(\eta\) 为血液的粘滞系数，\(r\) 为红细胞半径，\(v\) 为红细胞运动的速率。若某血样中半径为 \(r\) 的红细胞，由静止下沉直到匀速运动的速度为 \(v_m\)，红细胞密度为 \(\rho_1\)，血液的密度为 \(\rho_2\)。以下说法正确的是\paren[]
\begin{choices}
	\item 该红细胞先做匀加速运动，后做匀速运动
	\item 该红细胞的半径可表示为 \(r = \sqrt{\frac{9\eta v_m}{2g(\rho_1 - \rho_2)}}\)
	\item 若血样中红细胞的半径较小，则红细胞匀速运动的速度较大
	\item 若采用国际单位制中的基本单位来表示 \(v\) 的单位，则其单位为\qty{}{kg\cdot m^{-1}\cdot s^{-2}}
\end{choices}
\begin{answer}
    B\qquad 提示：
\end{answer}